\documentclass[10pt,a4paper]{amsart}
\usepackage[margin=19mm]{geometry}
\usepackage{amsmath,amssymb,mathtools}
\usepackage[scr=rsfs,cal=euler]{mathalfa}
\usepackage{xfrac}
\usepackage[unicode,hidelinks,bookmarks=false]{hyperref}
\newcommand{\ie}{i.e.\ }
\newcommand{\cf}{cf.\ }
\newcommand{\resp}{respectively\ }
\newcommand{\Subsection}{\subsection}
\providecommand{\nobreakdash}{\nobreak\hbox{-}}
\setlength{\emergencystretch}{2em}
\multlinegap0pt
\usepackage{amssymb}
\usepackage{xcolor}
\usepackage{tikz}
\usepackage{mathtools}
\usepackage[shortlabels]{enumitem}
\newtheorem{lemma}{Lemma}
\newcommand{\alphastar}{\alpha^\star}
\newcommand{\pr}[1]{\left(#1\right)}
\newcommand{\rzero}[0]{r_{0}}
\title{Sharp Risk Bounds for Early-Stopping in Gaussian \\ Linear Regression}
\author{Tobias Wegel, Gil Kur, Patrick Rebeschini}
\date{Source: arXiv:2503.03426v2 (CC BY 4.0)}
\begin{document}
\maketitle

\noindent\emph{Source and licence.} Sharp Risk Bounds for Early-Stopping in Gaussian \\ Linear Regression. Version arXiv:2503.03426v2 (CC BY 4.0), distributed by arXiv under the Creative Commons Attribution 4.0 International licence, http://creativecommons.org/licenses/by/4.0/, as recorded in the arXiv licence metadata for this exact version and shown on the abstract page https://arxiv.org/abs/2503.03426v2. Full licence notice: \texttt{licenses/arxiv-2503.03426-CC-BY-4.0.txt}.\par\smallskip

\noindent\textit{Editorial scope.} Counted unit: the article's lemma eq:critical-radius-scaling (source0.tex lines 854--860). The article's complete original proof of it is delivered with it (source0.tex lines 862--876, one \texttt{proof} environment of 15 lines). No mathematics is written, repaired or re-derived: the statement and the proof are the article's own lines, in the article's own notation, and no retained line is substituted or re-flowed.  Retained uncounted as the source-local context the proof needs: lines 306--309.  Nothing was omitted from any retained span, so the package carries no omission record.  No retained line contains a citation, so the wrapper carries no bibliography and no bibliography entry.  No retained line contains a figure environment, an image-inclusion command or drawing code, so no image is delivered and the package declares no asset dependency.  Editorial formatting only, in addition: an AMS \texttt{amsart} preamble that restates the article's own declarations and macros used by the retained lines, namely the environments \texttt{lemma} on separate counters, together with the standard packages those lines need.  The retained environments print in this document as Lemma 1 in order of appearance, whereas the article prints the same items with the numbering of its own full document.  The complete native source of the article as distributed in the arXiv e-print for this exact version is bundled unchanged as \texttt{sources/174254/source0.tex}.
\begin{equation}
\label{eq:localized-inequality}
    w\pr{(K-\alpha) \cap r B^d_2}\leq \frac{r^2}{2}
\end{equation}

\begin{lemma}
    For any convex body $K$ that contains $0$ in its interior, and any $\tau\geq 1$, it holds that
\begin{equation}
\label{eq:critical-radius-scaling}
    \rzero^2(\tau K):=\sup_{\alphastar\in \tau K}\rzero^2(\alphastar,\tau K)\leq \tau \cdot \sup_{\alphastar\in K} \rzero^2(\alphastar,K) =: \tau \cdot \rzero^2(K).
\end{equation}
\end{lemma}

\begin{proof}
    Note that, by definition, $\rzero(K)$ is bounded by any $r\geq 0$ that satisfies for all $\alphastar \in K$ (cf. \eqref{eq:localized-inequality})
\begin{equation*}
     w\pr{(K-\alphastar)\cap rB_2}\leq \frac{r^2}{2}.
\end{equation*}
Take any $\alphastar \in \tau K$ and define $R = \sqrt{\tau} \cdot \rzero(K)$. We have that
\begin{align*}
    w\pr{(\tau K-\alphastar)\cap RB_2} &= w\pr{\tau \pr{( K-\alphastar/\tau)\cap (R/\tau) B_2}}\\
    &=\tau \cdot w\pr{( K-\alphastar/\tau)\cap (R/\tau) B_2} \\
    &= \tau \cdot w\pr{( K-\alphastar/\tau)\cap (\rzero(K)/\sqrt{\tau}) B_2} \\
    &\leq \tau \cdot w\pr{( K-\alphastar/\tau)\cap \rzero(K) B_2} \\
    &\leq \tau \cdot \frac{\rzero^2(K)}{2} = \frac{R^2}{2}.
\end{align*}
Consequently, $\rzero^2(\alphastar,\tau K)\leq R^2 = \tau \rzero^2(K)$ for every $\alphastar\in \tau K$, and hence we have that $\sup_{\alphastar \in \tau K}\rzero^2(\alphastar, \tau K)\leq \tau \cdot \sup_{\alphastar \in K}\rzero^2(\alphastar, K)$, or in short, $\rzero^2(\tau K) \leq \tau \cdot \rzero^2(K)$.
\end{proof}

\end{document}
